\documentclass[11pt, letterpaper]{article}

% --- Idioma y codificación ---
\usepackage[spanish]{babel}
\usepackage[utf8]{inputenc}
\usepackage[T1]{fontenc}
\usepackage{lmodern}

% --- Matemáticas ---
\usepackage{amsmath, amssymb, amsthm}

% --- Formato ---
\usepackage[margin=2.5cm]{geometry}
\usepackage{parskip}
\usepackage{booktabs}
\usepackage{hyperref}
\hypersetup{colorlinks=true, linkcolor=blue!70!black, urlcolor=blue!70!black}

% --- Entornos de teoremas ---
\theoremstyle{definition}
\newtheorem{definition}{Definición}[section]
\theoremstyle{plain}
\newtheorem{theorem}{Teorema}[section]
\newtheorem{lemma}[theorem]{Lema}
\newtheorem{corollary}[theorem]{Corolario}
\theoremstyle{remark}
\newtheorem{remark}{Observación}[section]

\title{%
  T02 --- Demostración formal:\\[4pt]
  Correctitud de Bellman-Ford por inducción
}
\author{Curso Linux--C--Rust --- B15 Estructuras de datos en C}
\date{}

\begin{document}
\maketitle
\tableofcontents

% ===================================================================
\section{Modelo}
% ===================================================================

\begin{definition}
Grafo $G = (V, E)$ dirigido con $w: E \to \mathbb{R}$ (pesos
arbitrarios). Fuente $s \in V$, $|V| = n$, $|E| = m$.
$\delta(s,v)$ = peso del camino más corto ($-\infty$ si ciclo
negativo alcanzable, $\infty$ si inalcanzable).
$\delta_k(s,v)$ = peso del camino más corto con a lo sumo $k$
aristas. Bellman-Ford realiza $n - 1$ pasadas; en cada una relaja
todas las aristas.
\end{definition}

% ===================================================================
\section{Lema previo: caminos más cortos simples}
% ===================================================================

\begin{lemma}\label{lem:simple}
Si no existen ciclos negativos alcanzables desde $s$, entonces para
todo $v$ alcanzable existe un camino más corto simple con a lo sumo
$n - 1$ aristas.
\end{lemma}

\begin{proof}
Sea $p$ un camino más corto de $s$ a $v$. Si $p$ contiene un ciclo
$c$, su peso es $w(c) \geq 0$ (no hay ciclos negativos). Eliminando
$c$:
\[
  w(p') = w(p) - w(c) \leq w(p) = \delta(s,v)
\]
Como $p$ era más corto, $w(p') = w(p)$. Repitiendo hasta eliminar
todos los ciclos: camino simple. Un camino simple visita a lo sumo
$n$ vértices, usando a lo sumo $n - 1$ aristas.
\end{proof}

% ===================================================================
\section{Demostración 1: correctitud por inducción}
% ===================================================================

\begin{theorem}\label{thm:correctness}
Tras la pasada $i$ ($i = 1, \ldots, n{-}1$):
\[
  \mathrm{dist}[v] \leq \delta_i(s,v) \quad \forall\, v \in V
\]
En particular, tras $n - 1$ pasadas,
$\mathrm{dist}[v] = \delta(s,v)$ para todo~$v$ (sin ciclos
negativos).
\end{theorem}

\begin{proof}
Por inducción fuerte sobre $i$.

\textbf{Base ($i = 0$).} $\mathrm{dist}[s] = 0 = \delta_0(s,s)$.
Para $v \neq s$:
$\mathrm{dist}[v] = \infty = \delta_0(s,v)$. \checkmark

\textbf{Paso inductivo.} Supongamos que tras la pasada $i - 1$:
\[
  \mathrm{dist}[v] \leq \delta_{i-1}(s,v) \quad \forall\, v
  \tag{H.I.}
\]

Sea $v$ con $\delta_i(s,v) < \infty$. Sea
$p = s \leadsto u \to v$ un camino más corto con a lo sumo $i$
aristas. El subcamino $s \leadsto u$ tiene a lo sumo $i - 1$
aristas, y por subestructura óptima:
$w(s \leadsto u) = \delta_{i-1}(s,u)$.

Por H.I., tras la pasada $i - 1$:
\[
  \mathrm{dist}[u] \leq \delta_{i-1}(s,u) \tag{1}
\]

En la pasada $i$, se relaja $(u,v)$:
\[
  \mathrm{dist}[v] \leq \mathrm{dist}[u] + w(u,v) \tag{2}
\]

Combinando (1) y (2):
\[
  \mathrm{dist}[v] \leq \delta_{i-1}(s,u) + w(u,v)
  = \delta_i(s,v) \tag{3}
\]

Además, $\mathrm{dist}[v] \geq \delta(s,v)$ siempre (invariante de
cota superior, misma prueba que en Dijkstra).

\textbf{Tras $n - 1$ pasadas.} Por el Lema~\ref{lem:simple},
$\delta(s,v) = \delta_{n-1}(s,v)$:
\[
  \delta(s,v) \leq \mathrm{dist}[v]
  \leq \delta_{n-1}(s,v) = \delta(s,v)
  \quad\Longrightarrow\quad
  \mathrm{dist}[v] = \delta(s,v) \qedhere
\]
\end{proof}

% ===================================================================
\section{Demostración 2: detección de ciclos negativos}
% ===================================================================

\begin{theorem}\label{thm:negcycle}
Tras $n - 1$ pasadas, si existe $(u,v) \in E$ con
$\mathrm{dist}[u] + w(u,v) < \mathrm{dist}[v]$, entonces hay un
ciclo negativo alcanzable desde $s$.
\end{theorem}

\begin{proof}[Prueba por contraposición]
Si no hay ciclos negativos, por el Teorema~\ref{thm:correctness}:
$\mathrm{dist}[v] = \delta(s,v)$ para todo $v$. Para cualquier
$(u,v) \in E$, la desigualdad triangular da:
\[
  \delta(s,v) \leq \delta(s,u) + w(u,v)
\]

Sustituyendo:
$\mathrm{dist}[v] \leq \mathrm{dist}[u] + w(u,v)$. La condición
de relajación es falsa para toda arista. Contrapositiva: si alguna
arista es relajable, hay ciclo negativo.
\end{proof}

\begin{theorem}[Recíproco]\label{thm:negcycle-conv}
Si existe un ciclo negativo alcanzable desde $s$, la pasada
$n$-ésima detecta al menos una arista relajable.
\end{theorem}

\begin{proof}
Sea $c = v_0 \to v_1 \to \cdots \to v_k \to v_0$ un ciclo negativo
con $\sum_{j=0}^{k} w(v_j, v_{j+1}) < 0$.

Supongamos por contradicción que ninguna arista del ciclo es
relajable tras $n - 1$ pasadas:
\[
  \mathrm{dist}[v_{j+1}] \leq \mathrm{dist}[v_j] + w(v_j, v_{j+1})
  \quad \forall\, j = 0, \ldots, k
\]

Sumando sobre las $k + 1$ aristas:
\[
  \sum_{j=0}^{k} \mathrm{dist}[v_{j+1}]
  \leq \sum_{j=0}^{k} \mathrm{dist}[v_j]
  + \sum_{j=0}^{k} w(v_j, v_{j+1})
\]

Ambas sumas de $\mathrm{dist}$ son iguales (mismos vértices del
ciclo):
\[
  0 \leq \sum_{j=0}^{k} w(v_j, v_{j+1})
\]

Pero $\sum w(v_j, v_{j+1}) < 0$. Contradicción. Por tanto, al menos
una arista es relajable.
\end{proof}

% ===================================================================
\section{Demostración 3: complejidad
  \texorpdfstring{$\Theta(nm)$}{Theta(nm)}}
% ===================================================================

\begin{theorem}\label{thm:complexity}
Bellman-Ford ejecuta en $\Theta(nm)$ en el peor caso.
\end{theorem}

\begin{proof}
\emph{Cota superior.} Inicialización $O(n)$, $n - 1$ pasadas de
$m$ relajaciones, verificación $O(m)$:
\[
  T = O(n + (n{-}1)m + m) = O(nm)
\]

\emph{Cota inferior.} Camino
$s \to v_1 \to \cdots \to v_{n-1}$ con aristas en orden inverso en
la lista. En la pasada~1, solo $(s,v_1)$ relaja exitosamente (es la
última procesada). En la pasada~2, solo $(v_1,v_2)$. Se necesitan
exactamente $n - 1$ pasadas de $m$ aristas:
$T = \Omega(nm)$.

Combinando: $T = \Theta(nm)$.
\end{proof}

% ===================================================================
\section{Demostración 4: terminación temprana}
% ===================================================================

\begin{theorem}\label{thm:early}
Si una pasada completa no produce ninguna relajación exitosa, las
distancias son correctas.
\end{theorem}

\begin{proof}
Si ninguna arista $(u,v)$ satisface
$\mathrm{dist}[u] + w(u,v) < \mathrm{dist}[v]$, entonces:
\[
  \mathrm{dist}[v] \leq \mathrm{dist}[u] + w(u,v)
  \quad \forall\, (u,v) \in E
\]

Esto es la condición de la pasada $n$-ésima. Por el
Teorema~\ref{thm:negcycle}, no hay ciclos negativos, y por el
Teorema~\ref{thm:correctness},
$\mathrm{dist}[v] = \delta(s,v)$ para todo $v$. Pasadas adicionales
no cambiarían nada.
\end{proof}

\begin{remark}
Mejor caso: si el grafo no tiene pesos negativos y las aristas están
en ``buen'' orden (ej: BFS order), puede bastar una sola pasada:
$O(m)$.
\end{remark}

% ===================================================================
\section{Demostración 5: contaminación intra-pasada}
% ===================================================================

\begin{theorem}\label{thm:intrapass}
Tras exactamente $k$ pasadas, si se usa el array actualizado dentro
de la misma pasada (implementación estándar):
\[
  \mathrm{dist}^{(k)}[v] \leq \delta_k(s,v)
\]
La convergencia es al menos tan rápida como la versión pura (que usa
una copia del array anterior).
\end{theorem}

\begin{proof}
La desigualdad $\mathrm{dist}^{(k)}[v] \leq \delta_k(s,v)$ se
mantiene porque las actualizaciones intra-pasada solo pueden
\textbf{reducir} valores de $\mathrm{dist}$, propagando información
adicional. Nunca aumentan un $\mathrm{dist}$ ya correcto.

La versión pura satisface $\mathrm{dist}^{(k)}[v] = \delta_k(s,v)$
exactamente. La versión con contaminación satisface
$\mathrm{dist}^{(k)}[v] \leq \delta_k(s,v)$: puede converger más
rápido, nunca peor.
\end{proof}

\begin{corollary}
La implementación estándar (sin copia) nunca necesita más pasadas que
la versión pura, y frecuentemente menos.
\end{corollary}

% ===================================================================
\section{Resumen de resultados}
% ===================================================================

\begin{table}[h]
\centering
\begin{tabular}{llp{5.5cm}}
\toprule
\textbf{Resultado} & \textbf{Técnica} & \textbf{Clave} \\
\midrule
$\mathrm{dist}[v] = \delta(s,v)$ tras $n{-}1$ pasadas
  & Inducción sobre pasadas
  & Pasada $i$ corrige caminos de $\leq i$ aristas \\
Detección de ciclos negativos
  & Contrapositiva + contradicción
  & Sumar desigualdades en el ciclo: $0 \leq w(c) < 0$ \\
$\Theta(nm)$
  & Cota sup.\ + cota inf.
  & Peor caso: 1 arista corregida por pasada \\
Terminación temprana
  & Sin relajación $\Rightarrow$ desig.\ triangular
  & Equivale a verificación de pasada $n$ \\
Contaminación intra-pasada
  & $\mathrm{dist}^{(k)}[v] \leq \delta_k(s,v)$
  & Solo converge más rápido, nunca peor \\
\bottomrule
\end{tabular}
\end{table}

\end{document}
